\section{Research-program overview}

The program is organized as three primary neural aggregation levels plus the
benchmark and provenance infrastructure that connects them.

\begin{enumerate}
\item \textbf{Level I --- Representation formation.} How patch and
  foundation-model representations encode tissue, scanner, stain, acquisition
  workflow, site or center, region identity, and dataset artifacts. Principal
  research lines: PCam patch foundation; paired-scanner and paired-acquisition
  research; Paired-Acquisition Neural Factorization; SCORPION; multi-scanner
  canine SCC; scanner and center subspaces; centroid/QR, PCA, paired-linear,
  adversarial, and broken-pair controls; and synthetic identifiability and
  capacity-allocation studies.
\item \textbf{Level II --- Whole-slide aggregation.} How patch representations
  are aggregated into slide predictions. Principal research line: TransnnMIL and
  all of its implemented architectural families (canonical dual-branch,
  three-branch v2, branch-token fusion, hierarchical pooling, topology branch,
  adaptive pruning, graph caching).
\item \textbf{Level III --- Institutional aggregation.} How multi-institution
  updates, validation priorities, contribution signals, and safety constraints
  are aggregated. Principal research lines: PathologyFL; FAIR-WEIGHTS-H;
  institutional weighting; dominant-site stress testing; site imbalance;
  ordinal-grading shift; corruption and label-noise stress; detector behavior;
  client heterogeneity, dropout, and stragglers; communication accounting;
  secure aggregation; differential-privacy support; Byzantine or anomaly
  monitoring; and CAMELYON17 held-out-center and center-weighting studies.
\end{enumerate}

The three levels are connected by a shared evaluation and accountability
framework, and by shared frozen feature inputs (DINOv2, Phikon
\cite{filiot2023phikon}, ImageNet ResNet): scanner leakage, center leakage, descriptive task accessibility,
region retrieval, whole-slide aggregation, institutional influence, fairness
constraints, robustness, provenance, and claim validity are audited with the same
fail-closed discipline at every level. PCam provides the patch-level benchmark
foundation; PANDA provides the whole-slide benchmark; CAMELYON17 provides the
multi-center benchmark; and the provenance system binds every active claim to an
immutable artifact.

The unifying thesis is that computational-pathology neural systems aggregate
information at several levels --- patches into representations, patch instances
into whole-slide predictions, and institutional updates into collaborative
models --- and at each level nuisance variation, shortcut signals, optimization
failures, and institutional imbalance can be preserved or amplified. This program
develops the neural architectures, paired-data designs,
institutional-weighting protocols, stress tests, and provenance controls for
auditing and improving those aggregation stages.

\begin{equation}
\begin{aligned}
h_{r,s} &= E_{\omega}(x_{r,s}), \\
\mathcal{F}_{\mathrm{pair}}:\ \{h_{r,s}\}_{s\in\mathcal{S}_{r}}
&\longmapsto \left(z^{\mathrm{tissue}}_{r,s},z^{\mathrm{acq}}_{r,s}\right), \\
\mathcal{A}_{\mathrm{slide}}:\ \{(z_{k,j},p_{k,j})\}_{j\in\mathcal{B}_{k}}
&\longmapsto \widehat y_{k}, \\
\mathcal{A}_{\mathrm{site}}:\ \{(\Delta\theta_{i}^{(t)},w_{i}^{(t)})\}_{i\in\mathcal{I}_{t}}
&\longmapsto \theta^{(t+1)}, \\
\mathcal{P}_{\mathrm{PA\text{-}NF}}
&= \mathcal{A}_{\mathrm{site}}\circ\mathcal{A}_{\mathrm{slide}}\circ\mathcal{F}_{\mathrm{pair}}\circ E_{\omega}.
\end{aligned}
\label{eq:program}
\end{equation}

\noindent Here $E_{\omega}$ is a frozen or task-trained visual encoder, $\mathcal{F}_{\mathrm{pair}}$ is the paired-acquisition representation stage, $\mathcal{A}_{\mathrm{slide}}$ is whole-slide aggregation (TransnnMIL and matched MIL baselines), and $\mathcal{A}_{\mathrm{site}}$ is institutional aggregation (PathologyFL, FedAvg-family baselines, and dominance-aware weighting). Spatial coordinates $p_{k,j}$ are optional inputs to the slide stage for spatial and topology experiments.

\begin{figure}[!htbp]
\centering
\small
\caption{Publication-portfolio map. The flagship foundations manuscript frames
the program; focused papers A--D carry scoped evidence. Release order is
flagged by evidence maturity.}
\label{fig:portfolio}
\resizebox{\columnwidth}{!}{\begin{tabular}{lll}
\toprule
Component & Scope & Release readiness \\
\midrule
Flagship foundations manuscript & all three levels + audit & internal-review ready \\
Focused Paper A (PA-NF) & representation; negative increment; Layer-2 & corrected evidence ready \\
Focused Paper B (TransnnMIL) & whole-slide architecture & pending matched reruns \\
Focused Paper C (FL/FAIR-WEIGHTS-H) & institutional; infrastructure & after implementation completion \\
Focused Paper D (provenance) & audit system & portfolio decision \\
\bottomrule
\end{tabular}}
\end{figure}


\noindent\textbf{Evidence-maturity discipline.} Throughout, the manuscript keeps
four distinct categories separate: \emph{implemented architecture} (binds to
source and tests), \emph{corrected empirical evidence} (binds to immutable
result artifacts), \emph{negative or mixed results} (reported directly), and
\emph{pending controlled validation} (labeled pending). No component is implied
to have the same evidence maturity as any other.
